Let
\[
G=\int_{-1}^1K(u)S(u)S(u)^\top\,du.
\]
This matrix is positive definite because a polynomial whose squared integral against a kernel bounded below on a neighbourhood is zero must vanish identically. For \(h\le\min(t_0,1-t_0)\), the substitution \(a=t_0+hu\) and \(f_A\ge c\) give
\[
D_h
=
\int_{-1}^1K(u)S(u)S(u)^\top f_A(t_0+hu)\,du
\succeq cG.
\]
Thus \(D_h\) is invertible with uniformly bounded inverse. If \(q(t_0+hu)=\gamma_q^\top S(u)\) has degree at most \(p\), then
\[
\int W_h(a)q(a)f_A(a)\,da
=e_0^\top D_h^{-1}D_h\gamma_q=q(t_0).
\]
By Lemma~\texttt{lem:oracle-pseudo-outcome-identity},
\[
\vartheta_h=\int W_h(a)\theta_P(a)f_A(a)\,da.
\]
Let \(T_{t_0}\) be the Taylor polynomial of \(\theta_P\) at \(t_0\), of degree \(p\) for noninteger \(\alpha\) and degree \(p-1\) for integer \(\alpha=p\). Lemma~\texttt{lem:partial-mean-holder-transfer} and the frozen integer convention give
\[
\lvert\theta_P(a)-T_{t_0}(a)\rvert\le C\lvert a-t_0\rvert^\alpha.
\]
Moreover, the inverse bound and \(f_A\le1/c\) imply \(\int|W_h(a)|f_A(a)\,da\le C\). Polynomial reproduction therefore yields \(\lvert\vartheta_h-\theta_P(t_0)\rvert\le Ch^\alpha\). Only two-sided positivity of \(f_A\) was used.